\documentclass[10pt,a4paper]{amsart}
\usepackage[margin=19mm]{geometry}
\usepackage{amsmath,amssymb,mathtools}
\usepackage[scr=rsfs,cal=euler]{mathalfa}
\usepackage{xfrac}
\usepackage[unicode,hidelinks,bookmarks=false]{hyperref}
\newcommand{\ie}{i.e.\ }
\newcommand{\cf}{cf.\ }
\newcommand{\resp}{respectively\ }
\newcommand{\Subsection}{\subsection}
\providecommand{\nobreakdash}{\nobreak\hbox{-}}
\setlength{\emergencystretch}{2em}
\multlinegap0pt
\usepackage{bm}
\usepackage{bbm}
\usepackage{dsfont}
\usepackage{xspace}
\usepackage{cancel}
\usepackage{mathtools}
\usepackage{amsthm}
\makeatletter
\@ifundefined{proposition}{\newtheorem{proposition}{Proposition}}{}
\makeatother
\title[Very-Well-Behaved Epireflections for Categories of Models of]{Very-Well-Behaved Epireflections for Categories of Models of Sketches}
\author{Jo\~ao J. Xarez}
\date{Source: arXiv:2509.07241v1 (CC BY 4.0)}
\begin{document}
\maketitle

\noindent\emph{Source and licence.} Very-Well-Behaved Epireflections for Categories of Models of Sketches. Version arXiv:2509.07241v1 (CC BY 4.0), distributed under the Creative Commons Attribution 4.0 International licence, as recorded in the arXiv licence metadata for this exact version and shown on the abstract page https://arxiv.org/abs/2509.07241v1. Full rights notice: licenses/arxiv-2509.07241-CC-BY-4.0.txt.\par\smallskip

\noindent\textit{Editorial scope.} Counted unit: the Proposition whose source label is \texttt{proposition:idempotent endofunctor} in the article (source0.tex lines 129--133, source label \texttt{proposition:idempotent endofunctor}), delivered together with the article's own complete original proof of it (source0.tex lines 135--150, one \texttt{proof} environment of 16 lines). No mathematics is written, repaired or re-derived: statement and proof are the article's own text line for line. Required context retained uncounted: none. Every definition, display and label of those lines is retained unchanged. Omitted from this package and not redistributed: 1--128, 134--134, 151--1164. Nothing inside a retained excerpt is omitted or altered: every retained body is delivered exactly as the article's lines are written, so no omission receipt of any kind accompanies this package. Citations: the retained excerpts carry no citation command at all, so no bibliography entry of the article is transcribed and no reference is redistributed. Reference adaptation: none was needed and none was made; every reference inside the delivered unit resolves inside it, so no unresolved reference or citation is produced. Printed slips kept literal: no printed word, symbol, hypothesis or number of the counted statement or of its proof was altered. Layout and class: the article's own class is \texttt{amsart} with packages including amsmath, amsthm, amscd, amsfonts, graphicx, xy; the wrapper is the lane's pinned \texttt{amsart} editorial wrapper at 10pt on A4, which restates the theorem-like environments the retained text needs and adds the article's own macro definitions that the retained text needs (none), with extra wrapper packages (bm, bbm, dsfont, xspace, cancel, mathtools). The delivered numbering is the unit's own. Assets: no retained span contains a figure environment, a graphics-inclusion command, a PDF-page inclusion, a table or a TikZ picture; no image file is copied into this package, the package declares no asset dependency and no image file is redistributed with it. Licence: version arXiv:2509.07241v1 of this article is distributed under the Creative Commons Attribution 4.0 International licence, as recorded in the arXiv licence metadata for this exact version and shown on the abstract page https://arxiv.org/abs/2509.07241v1. A full rights notice is delivered at licenses/arxiv-2509.07241-CC-BY-4.0.txt. Characters: every retained excerpt is pure ASCII, so the delivered engine, XeLaTeX with its default Latin Modern text font, reports no missing character.
\begin{proposition}\label{proposition:idempotent endofunctor}

A well-pointed endofunctor $(\mathcal{I},\eta)$ on a category $\mathbb{C}$ is idempotent (i.e., $\eta\mathcal{I}$ is an isomorphism) if and only if $Fix(\mathcal{I},\eta)$ is a reflective full replete subcategory of $\mathbb{C}$ with reflection $\eta$, where $Fix(\mathcal{I},\eta)$ is determined by all the objects $M$($\in\mathbb{C}$) for which $\eta_M:M\rightarrow \mathcal{I}(M)$ is an isomorphism.

\end{proposition}

\begin{proof}
$Fix(\mathcal{I},\eta)$ is a replete subcategory of $\mathbb{C}$, since if $k:M\rightarrow N$ is an isomorphism from an object $M\in Fix(\mathcal{I},\eta)$, then $\eta_N=\mathcal{I}k\circ\eta_M\circ k^{-1}$ is also an isomorphism, assuring that $N\in Fix(\mathcal{I},\eta)$.

\noindent ($\Rightarrow$(only if))

If $(\mathcal{I},\eta)$ is idempotent, then any image $\mathcal{I}(C)$ belongs to $Fix(\mathcal{I},\eta)$, $C\in\mathbb{C}$ ($\eta \mathcal{I}$ is an isomorphism by hypothesis). Name $I=\mathcal{I}|_{Fix(\mathcal{I},\eta)}$ the co-restriction of the functor $\mathcal{I}$ to $Fix(\mathcal{I},\eta)$, and $H$ the full inclusion of $Fix(\mathcal{I},\eta)$ into $\mathbb{C}$. It is going to be shown that $I$ is left adjoint to $H$, with unit $\eta$.

Let $f:C\rightarrow M$ be any morphism whose codomain $M$ belongs to $Fix(\mathcal{I},\eta)$. Then, $\eta _M\circ f=\mathcal{I}f\circ \eta_C\Leftrightarrow f=(\eta^{-1}_M\circ \mathcal{I}f)\circ \eta_C$. It remains to check that $f'=\eta^{-1}_M\circ\mathcal{I}f$ is the unique morphism such that $f=f'\circ\eta_C$:
$$f''\circ\eta_C=f'\circ\eta_C\Rightarrow \mathcal{I}(f'')\circ\mathcal{I}(\eta_C)=\mathcal{I}(f')\circ\mathcal{I}(\eta_C)
$$
$$\Rightarrow \mathcal{I}(f'')\circ\eta_{\mathcal{I}(C)}=\mathcal{I}(f')\circ\eta_{\mathcal{I}(C)}\ (\mathcal{I}\eta=\eta\mathcal{I},\ by\ hypothesis)$$
$$\Rightarrow \eta_M\circ f''=\eta_M\circ f'\Rightarrow f''=f'\ (M\in Fix(\mathcal{I},\eta)).$$

\noindent ($\Leftarrow$(if))

If $\mathcal{I}(C)$ is the reflection of $C$, then $\mathcal{I}(C)\in Fix(\mathcal{I},\eta)$ and $\eta_{\mathcal{I}(C)}$ is an isomorphism. Hence, $\mathcal{I}\eta =\eta\mathcal{I}$ is an isomorphism.\end{proof}

\end{document}
